\documentclass[10pt,a4paper]{amsart}
\usepackage{amsmath,amssymb,amsthm,graphicx}
% For commutative diagrams we use Xy-pic
\usepackage[all]{xy}

% We use 2cell for 2-commutative diagrams.
\xyoption{2cell}
% Theorem environments.
%
\theoremstyle{plain}
\newtheorem{theorem}[subsection]{Theorem}
\newtheorem{proposition}[subsection]{Proposition}
\newtheorem{lemma}[subsection]{Lemma}

\theoremstyle{definition}
\newtheorem{definition}[subsection]{Definition}
\newtheorem{example}[subsection]{Example}
\newtheorem{exercise}[subsection]{Exercise}
\newtheorem{situation}[subsection]{Situation}

\theoremstyle{remark}
\newtheorem{remark}[subsection]{Remark}
\newtheorem{remarks}[subsection]{Remarks}

\numberwithin{equation}{subsection}

% Macros
%
\def\lim{\mathop{\mathrm{lim}}\nolimits}
\def\colim{\mathop{\mathrm{colim}}\nolimits}
\def\Spec{\mathop{\mathrm{Spec}}}
\def\Hom{\mathop{\mathrm{Hom}}\nolimits}
\def\Ext{\mathop{\mathrm{Ext}}\nolimits}
\def\SheafHom{\mathop{\mathcal{H}\!\mathit{om}}\nolimits}
\def\SheafExt{\mathop{\mathcal{E}\!\mathit{xt}}\nolimits}
\def\Sch{\mathit{Sch}}
\def\Mor{\mathop{\mathrm{Mor}}\nolimits}
\def\Ob{\mathop{\mathrm{Ob}}\nolimits}
\def\Sh{\mathop{\mathit{Sh}}\nolimits}
\def\NL{\mathop{N\!L}\nolimits}
\def\CH{\mathop{\mathrm{CH}}\nolimits}
\def\proetale{{pro\text{-}\acute{e}tale}}
\def\etale{{\acute{e}tale}}
\def\QCoh{\mathit{QCoh}}
\def\Ker{\mathop{\mathrm{Ker}}}
\def\Im{\mathop{\mathrm{Im}}}
\def\Coker{\mathop{\mathrm{Coker}}}
\def\Coim{\mathop{\mathrm{Coim}}}

% Boxtimes
%
\DeclareMathSymbol{\boxtimes}{\mathbin}{AMSa}{"02}

%
% Macros for moduli stacks/spaces
%
\def\QCohstack{\mathcal{QC}\!\mathit{oh}}
\def\Cohstack{\mathcal{C}\!\mathit{oh}}
\def\Spacesstack{\mathcal{S}\!\mathit{paces}}
\def\Quotfunctor{\mathrm{Quot}}
\def\Hilbfunctor{\mathrm{Hilb}}
\def\Curvesstack{\mathcal{C}\!\mathit{urves}}
\def\Polarizedstack{\mathcal{P}\!\mathit{olarized}}
\def\Complexesstack{\mathcal{C}\!\mathit{omplexes}}
% \Pic is the operator that assigns to X its picard group, usage \Pic(X)
% \Picardstack_{X/B} denotes the Picard stack of X over B
% \Picardfunctor_{X/B} denotes the Picard functor of X over B
\def\Pic{\mathop{\mathrm{Pic}}\nolimits}
\def\Picardstack{\mathcal{P}\!\mathit{ic}}
\def\Picardfunctor{\mathrm{Pic}}
\def\Deformationcategory{\mathcal{D}\!\mathit{ef}}
\usepackage{cleveref}
\crefname{originallemma}{original lemma}{original lemmas}
\Crefname{originallemma}{Original lemma}{Original lemmas}
\makeatletter
\newlabel{lemma-compare-higher-direct-image}{{075A}{source}{}{sourcepointer-1}{}}
\newlabel{lemma-compare-higher-direct-image@cref}{{[originallemma][0][]075A}{[1][0][]source}{}{}{}}
\makeatother
\crefname{originallemma}{original lemma}{original lemmas}
\Crefname{originallemma}{Original lemma}{Original lemmas}
\makeatletter
\newlabel{lemma-describe-pullback}{{09XN}{source}{}{sourcepointer-2}{}}
\newlabel{lemma-describe-pullback@cref}{{[originallemma][0][]09XN}{[1][0][]source}{}{}{}}
\makeatother
\crefname{originallemma}{original lemma}{original lemmas}
\Crefname{originallemma}{Original lemma}{Original lemmas}
\makeatletter
\newlabel{lemma-proper-base-change}{{09V5}{source}{}{sourcepointer-3}{}}
\newlabel{lemma-proper-base-change@cref}{{[originallemma][0][]09V5}{[1][0][]source}{}{}{}}
\makeatother
\crefname{originallemma}{original lemma}{original lemmas}
\Crefname{originallemma}{Original lemma}{Original lemmas}
\makeatletter
\newlabel{lemma-proper-base-change-f-star}{{0A3U}{source}{}{sourcepointer-4}{}}
\newlabel{lemma-proper-base-change-f-star@cref}{{[originallemma][0][]0A3U}{[1][0][]source}{}{}{}}
\makeatother
\crefname{originallemma}{original lemma}{original lemmas}
\Crefname{originallemma}{Original lemma}{Original lemmas}
\makeatletter
\newlabel{lemma-proper-base-change-mod-n}{{0F0C}{source}{}{sourcepointer-5}{}}
\newlabel{lemma-proper-base-change-mod-n@cref}{{[originallemma][0][]0F0C}{[1][0][]source}{}{}{}}
\makeatother
\crefname{originalproposition}{original proposition}{original propositions}
\Crefname{originalproposition}{Original proposition}{Original propositions}
\makeatletter
\newlabel{proposition-finite-higher-direct-image-zero}{{03QP}{source}{}{sourcepointer-6}{}}
\newlabel{proposition-finite-higher-direct-image-zero@cref}{{[originalproposition][0][]03QP}{[1][0][]source}{}{}{}}
\makeatother
\crefname{originalremark}{original remark}{original remarks}
\Crefname{originalremark}{Original remark}{Original remarks}
\makeatletter
\newlabel{sites-cohomology-remark-base-change}{{07A7}{source}{}{sourcepointer-7}{}}
\newlabel{sites-cohomology-remark-base-change@cref}{{[originalremark][0][]07A7}{[1][0][]source}{}{}{}}
\makeatother
\title{Small and big pushforwards agree for proper morphisms, for torsion coefficients and for Z/nZ coefficients}
\author{The Stacks Project Authors}
\date{stacks-project:etale-cohomology.tex:local-snapshot-20261011; GNU FDL 1.2 or later}
\begin{document}
\maketitle
\noindent\textit{Editorial note.} All delivered mathematical text is the author's own, copied from the pinned Stacks Project chapter file etale-cohomology.tex; no statement or proof step was rewritten, reordered or removed.\par\medskip
\noindent\textit{Editorial note.} The counted proof is the author's complete argument: the proof is written out in full for all four parts: the formal unwinding in (1), the pullback computation identifying the base change map in (2), the same argument with modulo n coefficients in (3), and the exactness reductions for finite morphisms in (4).\par\medskip
\noindent\textit{Editorial note.} The label lemma-compare-higher-direct-image-proper names the counted statement, and the proof environment immediately following it in the source is delivered verbatim as its proof.\par\medskip
\noindent\textit{Editorial note.} The statement and proof are copied from the pinned chapter file; the seven cited inputs are kept as external pointers with their canonical Stacks tags.\par\medskip
\noindent\textit{Editorial note.} The preamble is the chapter's own native preamble excerpt already used for this source: the Xy-pic setup (native preamble lines 14-18) and the native theorem declarations with the mathematical macros (native preamble lines 154-217).\par\medskip
\noindent\textit{Allowed standard prerequisites (editorial).}\begin{itemize}
\item The comparison morphisms pi\_S and the functors f\_small,*, f\_big,*. Same chapter, the morphism pi\_S : (Sch/S)\_etale -> S\_etale and the small and big pushforwards of a morphism of schemes f : T -> S. Supplies both sides of every comparison asserted in the four parts.
\item Base change of the comparison map along an object of (Sch/S)\_etale. Cohomology on Sites, the remark providing the canonical base change map pi\_S\textasciicircum{}\{-1\}Rf\_\{small,*\}K -> Rf\_\{big,*\}pi\_T\textasciicircum{}\{-1\}K. Supplies the map whose being an isomorphism is proved in parts (2) and (3).
\item Proper base change for torsion coefficients and modulo n coefficients. Same chapter, the proper base change theorem and its Z/nZ variant. Supplies the isomorphism of the pulled-back base change map in parts (2) and (3).
\item Exactness of the pushforward along a finite morphism. Same chapter, the vanishing of the higher direct images of a finite morphism. Supplies exactness of f\_\{small,*\} and, via the comparison lemma, that of f\_\{big,*\} in part (4).
\end{itemize}\medskip
\ifdefined\hypertarget\hypertarget{sourcepointer-1}{}\fi
\noindent\textit{Source reference.} Original lemma 075A: sources/stacks-etale-cohomology/pinned-original-source.tex, source lines 21196--21196. Prerequisite pointer invoked at proof line(s) 21393; it names a standard fact of the same chapter that the counted argument uses, and every step of the counted proof built on it is retained verbatim in the delivered body.\par\medskip
\ifdefined\hypertarget\hypertarget{sourcepointer-2}{}\fi
\noindent\textit{Source reference.} Original lemma 09XN: sources/stacks-etale-cohomology/pinned-original-source.tex, source lines 4960--4960. Prerequisite pointer invoked at proof line(s) 21354; it names a standard fact of the same chapter that the counted argument uses, and every step of the counted proof built on it is retained verbatim in the delivered body.\par\medskip
\ifdefined\hypertarget\hypertarget{sourcepointer-3}{}\fi
\noindent\textit{Source reference.} Original lemma 09V5: sources/stacks-etale-cohomology/pinned-original-source.tex, source lines 18708--18708. Prerequisite pointer invoked at proof line(s) 21395; it names a standard fact of the same chapter that the counted argument uses, and every step of the counted proof built on it is retained verbatim in the delivered body.\par\medskip
\ifdefined\hypertarget\hypertarget{sourcepointer-4}{}\fi
\noindent\textit{Source reference.} Original lemma 0A3U: sources/stacks-etale-cohomology/pinned-original-source.tex, source lines 18282--18282. Prerequisite pointer invoked at proof line(s) 21362; it names a standard fact of the same chapter that the counted argument uses, and every step of the counted proof built on it is retained verbatim in the delivered body.\par\medskip
\ifdefined\hypertarget\hypertarget{sourcepointer-5}{}\fi
\noindent\textit{Source reference.} Original lemma 0F0C: sources/stacks-etale-cohomology/pinned-original-source.tex, source lines 18887--18887. Prerequisite pointer invoked at proof line(s) 21400; it names a standard fact of the same chapter that the counted argument uses, and every step of the counted proof built on it is retained verbatim in the delivered body.\par\medskip
\ifdefined\hypertarget\hypertarget{sourcepointer-6}{}\fi
\noindent\textit{Source reference.} Original proposition 03QP: sources/stacks-etale-cohomology/pinned-original-source.tex, source lines 8245--8245. Prerequisite pointer invoked at proof line(s) 21406; it names a standard fact of the same chapter that the counted argument uses, and every step of the counted proof built on it is retained verbatim in the delivered body.\par\medskip
\ifdefined\hypertarget\hypertarget{sourcepointer-7}{}\fi
\noindent\textit{Source reference.} Original remark 07A7: sources/stacks-chow/tags.txt, source lines 9297--9297. Prerequisite pointer invoked at proof line(s) 21374; it names a standard fact of a sibling chapter that the counted argument uses, and every step of the counted proof built on it is retained verbatim in the delivered body.\par\medskip
\begin{lemma}
\label{lemma-compare-higher-direct-image-proper}
Let $f : T \to S$ be a proper morphism of schemes. Then we have
\begin{enumerate}
\item $\pi_S^{-1} \circ f_{small, *} = f_{big, *} \circ \pi_T^{-1}$
as functors $\Sh(T_\etale) \to \Sh((\Sch/S)_\etale)$,
\item $\pi_S^{-1}Rf_{small, *}K = Rf_{big, *}\pi_T^{-1}K$
for $K$ in $D^+(T_\etale)$ whose cohomology sheaves are torsion,
\item $\pi_S^{-1}Rf_{small, *}K = Rf_{big, *}\pi_T^{-1}K$
for $K$ in $D(T_\etale, \mathbf{Z}/n\mathbf{Z})$, and
\item $\pi_S^{-1}Rf_{small, *}K = Rf_{big, *}\pi_T^{-1}K$
for all $K$ in $D(T_\etale)$ if $f$ is finite.
\end{enumerate}
\end{lemma}

\begin{proof}
Proof of (1). Let $\mathcal{F}$ be a sheaf on $T_\etale$.
Let $g : S' \to S$ be an object of $(\Sch/S)_\etale$. Consider the
fibre product
$$
\xymatrix{
T' \ar[r]_{f'} \ar[d]_{g'} & S' \ar[d]^g \\
T \ar[r]^f & S
}
$$
Then we have
$$
(f_{big, *}\pi_T^{-1}\mathcal{F})(S') =
(\pi_T^{-1}\mathcal{F})(T') =
((g'_{small})^{-1}\mathcal{F})(T')  =
(f'_{small, *}(g'_{small})^{-1}\mathcal{F})(S')
$$
the second equality by Lemma \ref{lemma-describe-pullback}.
On the other hand
$$
(\pi_S^{-1}f_{small, *}\mathcal{F})(S') =
(g_{small}^{-1}f_{small, *}\mathcal{F})(S')
$$
again by Lemma \ref{lemma-describe-pullback}.
Hence by proper base change for sheaves of sets
(Lemma \ref{lemma-proper-base-change-f-star})
we conclude the two sets are canonically isomorphic.
The isomorphism is compatible with restriction mappings
and defines an isomorphism
$\pi_S^{-1}f_{small, *}\mathcal{F} = f_{big, *}\pi_T^{-1}\mathcal{F}$.
Thus an isomorphism of functors
$\pi_S^{-1} \circ f_{small, *} = f_{big, *} \circ \pi_T^{-1}$.

\medskip\noindent
Proof of (2). There is a canonical base change map
$\pi_S^{-1}Rf_{small, *}K \to Rf_{big, *}\pi_T^{-1}K$
for any $K$ in $D(T_\etale)$, see
Cohomology on Sites, Remark \ref{sites-cohomology-remark-base-change}.
To prove it is an isomorphism, it suffices to prove the pull back of
the base change map by $i_g : \Sh(S'_\etale) \to \Sh((\Sch/S)_\etale)$
is an isomorphism for any object $g : S' \to S$ of $(\Sch/S)_\etale$.
Let $T', g', f'$ be as in the previous paragraph.
The pullback of the base change map is
\begin{align*}
g_{small}^{-1}Rf_{small, *}K
& =
i_g^{-1}\pi_S^{-1}Rf_{small, *}K \\
& \to
i_g^{-1}Rf_{big, *}\pi_T^{-1}K \\
& =
Rf'_{small, *}(i_{g'}^{-1}\pi_T^{-1}K) \\
& =
Rf'_{small, *}((g'_{small})^{-1}K)
\end{align*}
where we have used $\pi_S \circ i_g = g_{small}$,
$\pi_T \circ i_{g'} = g'_{small}$, and
Lemma \ref{lemma-compare-higher-direct-image}.
This map is an isomorphism by the proper base change theorem
(Lemma \ref{lemma-proper-base-change}) provided $K$ is bounded
below and the cohomology sheaves of $K$ are torsion.

\medskip\noindent
The proof of part (3) is the same as the proof of part (2), except
we use Lemma \ref{lemma-proper-base-change-mod-n}
instead of Lemma \ref{lemma-proper-base-change}.

\medskip\noindent
Proof of (4). If $f$ is finite, then the functors
$f_{small, *}$ and $f_{big, *}$ are exact. This follows
from Proposition \ref{proposition-finite-higher-direct-image-zero}
for $f_{small}$. Since any base change $f'$ of $f$ is finite too,
we conclude from Lemma \ref{lemma-compare-higher-direct-image} part (3)
that $f_{big, *}$ is exact too (as the higher derived functors are zero).
Thus this case follows from part (1).
\end{proof}

\end{document}
